% !TEX encoding = UTF-8 Unicode

% Beispiel für ein LaTeX-Dokument im Format "seminarvorlage"
\documentclass[ngerman]{seminarvorlage}
% ngerman = Deutsch in neuer Rechtschreibung, alternativ english
\usepackage{babel} % automatische Sprachunterstützung

\usepackage[utf8]{inputenc} % Kodierung der Non-ASCII-Zeichen
\usepackage[T1]{fontenc} % Moderne Fonts, Trennung von Wörtern mit Umlauten
\usepackage{cleveref} % für bequeme Referenzen, siehe \cref unten

\usepackage{embrac}% upright brackets in emphasised text, () and [], empfohlen

%\usepackage{biblatex}% besser als bibtex, aber dann biber statt bibtex benutzen

\begin{document}

% Unbedingt angeben: Titel, Autoren
% Freiwillig: Adresse, E-Mail
\title{Evaluation von KI-Funktionen in modernen Entwicklungsumgebungen}
\numberofauthors{2}
\author{
  \alignauthor Konstantin Fastovski \\
    \email{konstantin@fastovski.org}
  \alignauthor Silas Rißler\\
  \email{}
}

\maketitle% Titelangaben produzieren, aber kein Inhaltsverzeichnis, \tableofcontents funktioniert nicht!
\newpage

\abstract{Dieser Artikel untersucht moderne Entwicklungsumgebungen, Text-Editoren und Coding Agenten aus Sicht ihrer Nutzer*innen. Hauptbestandteil sind die angebotenen Werkzeuge zur Einbindung von künstlicher Intelligenz (KI) in den Workflow von Software-Entwicklern. Diese Werkzeuge werden auf ihre Nutzerfreundlichkeit, Tranzparenz sowie Qualität untersucht. Daraus wird identifiziert, welche Design-Prinzipien sich bei Produkten durchsetzen und weshalb. Gleichzeitig werden auch die Fehler von unbeliebten Umsetzungen klar.}

\keywords{AI, KI, Software Engineering, Software-Entwicklung}

% ca. 1 Seite
\section{Einleitung}

% ca. 1 Seite
\section{Wie benutzen Entwickler*innen ihre IDE?}

% ca. 2 Seiten, Katalog der Bewertungskriterien
\section{Kriterien-Katalog}

% ca. 2 Seiten
\section{Vertrauens- und Tranzparenzentscheidungen}

% ca. 2 Seiten
\section{Feature und Designmuster}

% ca. 1 Seite
\section{Fazit}

% Figure example:
% \begin{figure}[htbp]
% \begin{center}
% \unitlength8mm % hier skalieren, falls gewünscht
% \begin{picture}(4,6)
% %außen
% \put(0,0){\line(1,0){4}}
% \put(4,0){\line(0,1){3}}
% \put(4,3){\line(-1,1){2}}
% \put(2,5){\line(-1,-1){2}}
% \put(0,3){\line(0,-1){3}}
% %innen
% \put(4,3){\line(-1,0){4}}
% \end{picture}
% \end{center}
% \caption{Eine typische Gummibärenfabrik
%          siehe auch~\protect\cite{Ivory2001}.}
% \label{fabrik}
% \end{figure}

% Table example:
% \begin{table}[htbp]
% \begin{center}
% \begin{tabular}{|c|c|c|}
% \hline
% Material & Tierart & essbar\\
% \hline
% Gummi & Bär & ja\\
% Lakritz & Schnecke & geht so\\
% \hline
% \end{tabular}
% \end{center}
% \caption{Das gibt es alles.}
% \label{tttabelle}
% \end{table}

\bibliographystyle{abbrv}
\bibliography{literatur} % Daten aus der Datei literatur.bib verwenden.

\end{document}
